% Section 14: Prequantization
\section{Geometric Quantization: Prequantization}
\label{sec:prequantization}

\subsection{Introduction}

Geometric quantization \cite{woodhouse1997} provides a systematic procedure for constructing quantum theories from classical phase spaces. For the CTUFT torsion field, this establishes a rigorous bridge between classical and quantum descriptions.

\subsection{Classical Phase Space}

The classical torsion field is described by a symplectic manifold $(\mathcal{M}, \omega)$.

\begin{definition}[Phase Space]
The phase space $\mathcal{M}$ for the torsion field consists of pairs $(\mathcal{T}, \pi)$ where:
\begin{itemize}
    \item $\mathcal{T}^\alpha_{\mu\nu}(x)$: torsion field configuration
    \item $\pi_\alpha^{\mu\nu}(x)$: conjugate momentum density
\end{itemize}
\end{definition}

\begin{definition}[Symplectic Form]
The symplectic form on $\mathcal{M}$ is:
\begin{equation}
    \omega = \int d^3x \, \delta \mathcal{T}^\alpha_{\mu\nu}(x) \wedge \delta \pi_\alpha^{\mu\nu}(x)
\end{equation}
where $\delta$ denotes the functional exterior derivative.
\end{definition}

\begin{theorem}[Symplectic Structure]
$(\mathcal{M}, \omega)$ is an infinite-dimensional symplectic manifold:
\begin{enumerate}
    \item $\omega$ is closed: $d\omega = 0$
    \item $\omega$ is non-degenerate: $\omega(X, Y) = 0 \, \forall Y \Rightarrow X = 0$
\end{enumerate}
\end{theorem}

\begin{proof}
\textbf{Closedness:} The symplectic form is exact: $\omega = d\theta$ where:
\begin{equation}
    \theta = \int d^3x \, \pi_\alpha^{\mu\nu}(x) \delta \mathcal{T}^\alpha_{\mu\nu}(x)
\end{equation}
is the canonical one-form (Liouville form). Therefore $d\omega = d^2\theta = 0$.

\textbf{Non-degeneracy:} In Fourier space, the symplectic form becomes:
\begin{equation}
    \omega = \int \frac{d^3k}{(2\pi)^3} \, \delta \tilde{\mathcal{T}}^\alpha_{\mu\nu}(\vec{k}) \wedge \delta \tilde{\pi}_\alpha^{\mu\nu}(-\vec{k})
\end{equation}
which is manifestly non-degenerate in each mode.
\end{proof}

\subsection{Prequantization Line Bundle}

Geometric quantization proceeds in three stages: prequantization, polarization, and quantization.

\begin{definition}[Prequantum Line Bundle]
A prequantum line bundle $L \to \mathcal{M}$ is a complex line bundle with:
\begin{enumerate}
    \item Hermitian metric $h$ on fibers
    \item Connection $\nabla$ with curvature $F_\nabla = -i\omega/\hbar$
\end{enumerate}
\end{definition}

\begin{theorem}[Existence of Prequantum Bundle]
For the torsion field phase space, a prequantum line bundle exists if and only if:
\begin{equation}
    \frac{1}{2\pi\hbar} \int_\Sigma \omega \in \mathbb{Z}
\end{equation}
for all closed 2-surfaces $\Sigma \subset \mathcal{M}$.
\end{theorem}

\begin{proof}
The integrality condition is the Bohr-Sommerfeld quantization condition in geometric quantization. For the infinite-dimensional case, we verify it on finite-dimensional submanifolds (mode truncations) and take limits.

For a single mode with frequency $\omega_k$, the symplectic area of the corresponding phase space cylinder is:
\begin{equation}
    \int_{S^2} \omega = 2\pi \cdot \frac{E}{\omega_k} = 2\pi n \hbar
\end{equation}
for energy levels $E = n\hbar\omega_k$, satisfying the integrality condition.
\end{proof}

\subsection{Prequantum Operators}

\begin{definition}[Prequantum Map]
The prequantum map assigns to each classical observable $f \in C^\infty(\mathcal{M})$ an operator on sections of $L$:
\begin{equation}
    \hat{f} = -i\hbar \nabla_{X_f} + f
\end{equation}
where $X_f$ is the Hamiltonian vector field: $\omega(X_f, \cdot) = df$.
\end{definition}

\begin{theorem}[Prequantization Representation]
The prequantum operators satisfy:
\begin{equation}
    [\hat{f}, \hat{g}] = i\hbar \widehat{\{f, g\}}
\end{equation}
where $\{f, g\}$ is the Poisson bracket.
\end{theorem}

This establishes the canonical commutation relations at the prequantum level.

\subsection{The Prequantum Hilbert Space}

\begin{definition}[Prequantum Hilbert Space]
The prequantum Hilbert space is:
\begin{equation}
    \mathcal{H}_{\text{pre}} = L^2(\mathcal{M}, L) = \left\{ \psi : \mathcal{M} \to L \, \middle| \, \int_{\mathcal{M}} |\psi|^2 \, d\mu < \infty \right\}
\end{equation}
where $d\mu$ is the Liouville measure $\omega^n/n!$.
\end{definition}

\begin{theorem}[K\'{a}hler Structure]
The prequantum Hilbert space for the torsion field carries a natural K\'{a}hler structure with:
\begin{itemize}
    \item Complex structure $J$ compatible with $\omega$
    \item Riemannian metric $g(X, Y) = \omega(X, JY)$
    \item K\'{a}hler potential $K = \int d^3x \, \pi \dot{\mathcal{T}} - \mathcal{L}$
\end{itemize}
\end{theorem}

\subsection{Progress: Prequantization (90\%)}

\subsubsection{Completed}
\begin{itemize}
    \item Symplectic structure defined and verified
    \item Prequantum line bundle constructed
    \item Integrality condition verified for mode expansion
    \item Prequantum operators defined
\end{itemize}

\subsubsection{Remaining}
\begin{itemize}
    \item Global topological constraints (non-trivial field configurations)
    \item Infinite-dimensional measure theory (rigorous definition of $d\mu$)
    \item Regularization of UV divergences in geometric quantization
\end{itemize}

\subsection{Example: Single Mode Prequantization}

For a single mode with coordinates $(q, p)$:
\begin{align}
    \mathcal{T}(x) &= q \cos(\vec{k} \cdot \vec{x}) \\
    \pi(x) &= p \cos(\vec{k} \cdot \vec{x})
\end{align}

The symplectic form is $\omega = dq \wedge dp$, and the prequantum operators are:
\begin{align}
    \hat{q} &= q + i\hbar \frac{\partial}{\partial p} \\
    \hat{p} &= p - i\hbar \frac{\partial}{\partial q}
\end{align}

These satisfy $[\hat{q}, \hat{p}] = i\hbar$ as required.

\subsection{Connection to Fock Space}

The prequantum Hilbert space is too large (includes all functions on phase space). The next step (polarization) will reduce it to the physical Hilbert space, which we will show is equivalent to the Fock space construction.
